%------------------------------------------------------------------------------%
% \chapter{Plano Cartesiano e Equações}
%------------------------------------------------------------------------------%

%------------------------------------------------------------------------------%
\begin{exercicios}{Equação da Circunferência}
    
    \begin{exercicio}
        Determine a equação da circunferência que possui centro em $C(3,6)$ e raio 4.
        \begin{answers}
            \tabto{2em} $(x-3)^2 + (x-6)^2 = 16$
        \end{answers}
    \end{exercicio}
    
    \begin{exercicio}
        Determine a equação da circunferência com raio $2$ e centro em $(-1,3)$ e esboce a circunferência.
    \end{exercicio}

    \begin{exercicio}
        Verifique se os pontos $(6,1)$, $(4,5)$, $(0,5)$, $(5,1)$ e $(-2,-3)$ estão dentro, 
        fora ou sobre a circunferência de raio $5$ e centro $(1,1)$. 
        Esboce a circunferência e os pontos.
    \end{exercicio}
        
    \begin{exercicio}
        Qual a equação da circunferência de raio~$4$ cujo centro é a intersecção das retas 
        \[x -2y +1 = 0 \quad \text{e} \quad x + y - 2 = 0\]
    \end{exercicio}
        
    \begin{exercicio}
        Determine as equações do círculo que satisfazem as condições dadas.
        \begin{alphalist}
            \item Raio $5$ e centro $(2, \, -3)$
            \item Raio $3$ e centro $(-2, \, -4)$
            \item Raio $5$ e centro na origem
            \item Centro $(2, \, -3)$ e passando por $(5, \, 2)$
            \item Centro $(-a, \, a)$ e raio $2a$
        \end{alphalist}
        \begin{answers}
            \begin{mathlist}[2]
                \item (x-2)^2 + (y+3)^2 = 25
                \item (x+2)^2 + (y+4)^2 = 9
                \item x^2 + y^2 = 25
                \item (x-2)^2 + (y+3)^2 = 34
                \item (x+a)^2 + (y-a)^2 = 4a^2 
            \end{mathlist}
        \end{answers}
    \end{exercicio}
    
    \begin{exercicio}
        Determine os coeficientes $a$, $b$ e $c$ da equação do círculo, 
        $ x^2 + y^2 + ax + by + c = 0$, que passa pelos pontos
        $P_1 = (-2,7)$, $P_2 = (-4,5)$ e $P_3 = (4,-3)$.
    \end{exercicio}
    
    \begin{exercicio}
        Construa a equação da circunferência inscrita no quadrado de vértices $(1,1)$, $(3,1)$, $(3,3)$~e~$(1,3)$.
    \end{exercicio}
        
    \begin{exercicio}
        O centro de uma circunferência é determinado pelo ponto médio do segmento 
        $PQ$, sendo $P(4,6)$ e $Q(2,10)$. Considerando que o raio dessa circunferência é 7, 
        determine sua equação.
        \begin{answers}
            \tabto{2em} $(x-3)^2+(y-8)^2=49$
        \end{answers}
    \end{exercicio}
        
    \begin{exercicio}
        O ponto $P(3,b)$ pertence à circunferência de centro no ponto $C(0,3)$ e raio 5. Calcule valor da coordenada $b$. 
        \begin{answers}
            \tabto{2em} $b=-1$ ou $b=7$
        \end{answers}
    \end{exercicio}
    
    \begin{exercicio}
        Determine a equação da circunferência com centro no ponto $C(2,1)$ e que passa pelo ponto $A(1,1)$. 
        \begin{answers}
            \tabto{2em} $(x-2)^2+(y-1)^2=1$
        \end{answers}
    \end{exercicio}

  \begin{exercicio}
      Determine a equação da circunferência de centro $C$ e raio $r$
      \begin{alphalist}
          \item $C(0,0)$                         \tabto{3cm} $r=3$
          \item $C(2,0)$                         \tabto{3cm} $r=4$
          \item $C(-1,-2)$                       \tabto{3cm} $r=5$
          \item $C(2,4)$                         \tabto{3cm} $r=1$
          \item $C(0,-3)$                        \tabto{3cm} $r=2$
          \item $C(\sfrac{1}{2},\,\sfrac{3}{2})$ \tabto{3cm} $r=4$
        \end{alphalist}
        \begin{answers}
            \begin{mathlist}[2]
              \item x ^2 + y^2 = 9 
              \item (x-2) ^2 + y^2 = 16
              \item (x+1) ^2 + (y+2)^2 = 25
              \item (x-2) ^2 + (y-4)^2 = 1 
              \item x ^2 + (y+3)^2 = 4 
              \item (x-\sfrac{1}{2})^2 + (y-\sfrac{3}{2})^2 = 16
            \end{mathlist}
        \end{answers}
    \end{exercicio}
        
    \begin{exercicio}
        Qual é a equação da circunferência de centro $C(1,2)$ que passa por $P(5,5)$?
    \end{exercicio}
    
    \begin{exercicio}
        Determine o centro e o raio das seguintes circunferências
        \begin{mathlist}
            \item (x-2)^2 + (y+1)^2 = 36
            \item (x+2)^2 + y^2 - 9 = 0
            \item y^2 = 16 - x^2
            \item x^2 + y^2 +8x + 6y = 0
            \item x^2 + y^2 -6x + 4y - 12 = 0
        \end{mathlist}
        \begin{answers}
            \begin{mathlist}[2]
                \item C( 2,-1) \qquad r=6
                \item C(-2, 0) \qquad r=3
                \item C( 0, 0) \qquad r=4
                \item C(-4,-3) \qquad r=5
                \item C( 3,-2) \qquad r=5
            \end{mathlist}
        \end{answers}
    \end{exercicio}
        
    \begin{exercicio}
        Esboce as regiões determinadas pelas inequações
        \begin{mathlist}[2]
            \item x^2 + y^2 \leq 9
            \item x^2 + y^2 \geq 4
            \item (x-1)^2 + (1+y)^2 < 9
            \item 1 < (2+x)^2 + y^2 < 4
        \end{mathlist}        
    \end{exercicio}
    
\end{exercicios}

%------------------------------------------------------------------------------%
